\documentclass[11pt,a4paper]{article}
\usepackage[margin=23mm]{geometry}
\usepackage{fontspec}\setmainfont{Latin Modern Roman}
\usepackage{amsmath,amssymb,xurl,hyperref}
\setlength{\parindent}{0pt}\setlength{\parskip}{6pt}
\newcommand{\code}[1]{\nolinkurl{#1}}
\title{S05 / NORMALIZER-005: erratum końców przedziałów}
\author{Notatka robocza dla Niirmata}\date{10 października 2026}
\begin{document}\maketitle
\textbf{Status: jawna korekta błędu implementacji rachunku.}
Przy porównaniu precyzji MARGINAL-006 wykryto, że zapisane jako dokładne
końce przedziału dla tego samego wymiernego budżetu nie przecinały się.
Przyczyną była konwersja \code{QQ(endpoint)} z typu MPFR RealNumber.
W Sage nie musi ona zachowywać dokładnej wartości binarnej: może wybrać
pobliską liczbę wymierną. Dla górnego końca balls512 liczby $1/3$
zwraca $1/3$, ściśle mniej niż dokładny binarny górny koniec.

W oryginalnym \code{FIBER_NORMALIZER_005.sage} użyto tej konwersji
zarówno podczas składania nowych przedziałów theta, jak i serializacji
pól \code{lower_exact}, \code{upper_exact}. Nie jest to zatem wyłącznie
błąd kosmetycznego wydruku. Pierwotne certyfikaty liczbowe 005 nie są
autorytatywnym świadectwem skierowanego rachunku i zostają zastąpione
poniższą korektą. Nie nadpisano ich bajtów ani historycznych manifestów.

\textbf{Naprawa.} Każdy koniec MPFR zamieniamy przez
\code{endpoint.exact_rational()}. Funkcja składania przedziałów przyjmuje
wyłącznie już dokładne elementy $\mathbb Z$ lub $\mathbb Q$;
nie dokonuje niejawnej konwersji z liczb rzeczywistych.
W checkerze dodano kontrolę ujemną starej konwersji oraz kontrolę,
że zapisane końce binarne mają mianowniki będące potęgami 2.

\textbf{Nowe źródło:} \code{FIBER_NORMALIZER_005_CORRECTED.sage}.
Nowe świadectwa to \code{NORMALIZER_005_CORRECTED_CERTIFICATE.json}
oraz \code{NORMALIZER_005_CORRECTED_PRECISION768.json}.
Ponowne uruchomienia 512/768 bitów zachowują wszystkie zadeklarowane
granice 005: symboliczną szerokość względną reszt i odwrotności
$<2^{-356}$; dla h=1,c=0 szerokości pełnego normalizatora i odwrotności
$<2^{-180}$, TV pełnej odpowiedzi zerowego aliasu $<2^{-180}$ oraz
moment sekwencyjny przed normą $>2^{1536}$.
Rachunek MARGINAL-006 został również naprawiony i otrzymuje wyłącznie
skorygowany normalizator jako wejście.

Dowód symboliczny oparty na całkowitych końcach Taylora i faktoryzacji
nie zmienił się. Zmieniła się implementacja numerycznych obudów i ich
wiążące świadectwa. Dawny pozytywny test przecinania się 50 przedziałów
005 nie mógł sam wykryć niepoprawnej semantyki konwersji.

Historia obejmuje nieudany \code{crosscheck_001}, źródła przed naprawą,
probe API i surowe logi. Receipty MARGINAL-006 podają ich piny i nowe
uruchomienia. Wyniki nadal mają status \textbf{textual}, bez nowego kernela,
niezależnego review ani twierdzenia bezpieczeństwa.
\end{document}
